% ЗГЕНЕРОВАНО scripts/build_format_v2_chapter.py — не редагувати вручну.
% Джерело: books/ashtavakra-gita/poetychnyi/rozdil-11.md

\poemverse{11.1}{Видозміна буття й небуття — \\ за власною природою. \\ Хто має таке переконання, \\ той без видозмін, \\ з минулими турбаціями, \\ втихає легко.}
\poemverse{11.2}{Творець усього — Ішвара, \\ і немає тут іншого. \\ Хто має таке переконання, \\ той із розчиненими всередині сподіваннями, \\ втихлий, \\ ніде не застрягає.}
\poemverse{11.3}{Нещастя й достаток у свій час — \\ саме від долі. \\ Хто має таке переконання, \\ той вдоволений, \\ із чуттями в порядку, \\ повсякчас не бажає й не сумує.}
\poemverse{11.4}{Щастя й страждання, \\ народження й смерть — \\ саме від долі. \\ Хто має таке переконання, \\ той, не бачачи нічого досяжного, \\ без натуги, навіть діючи, \\ не забруднюється.}
\poemverse{11.5}{Страждання народжується з думання, \\ не інакше. \\ Хто має таке переконання, \\ той, позбавлений його, \\ щасливий, втихлий, \\ скрізь із опалою спрагою.}
\poemverse{11.6}{Я не тіло, і тіло не моє; \\ я — усвідомлення. \\ Хто має таке переконання, \\ той, наче досягши окремішності, \\ не памʼятає незробленого як зробленого.}
\poemverse{11.7}{Від Брахми до стеблини — \\ це саме я. \\ Хто має таке переконання, \\ той без домислювання, \\ чистий, втихлий, \\ вдоволений і досягнутим, і недосягнутим.}
\poemverse{11.8}{Цей світ, \\ різноманітно дивовижний, — \\ аніщо. \\ Хто має таке переконання, \\ той без усталених схильностей, \\ саме лише сяйнування, \\ втихає, наче аніщо.}
